\chapter{The Bride of Nine Knots}

\section*{Told by}
Lady Ysoria, a Vhasian widow who lives in a red tower and receives visitors only at dusk. Her left hand has six fingers. She says the extra one was a wedding gift. She did not say who gave it. She also did not say whether she was the bride, the rival, or the priest. Some questions are not asked in the red tower. The wine is older than the house. The silences are older than the wine.

\section*{The Tale}

A bride tied her dowry rope with nine knots. One for each year she waited for her betrothed while he fought wars in the southern passes. The priest said, \textit{``Eight is custom. Nine is a door. Do not tie what you cannot untie.''}

She tied it anyway. The ninth knot she tied with her teeth, because her hands were trembling.

The morning of the wedding, her rival — who had also loved the groom, who had also waited, whose own dowry rope had only seven knots because her family was poorer and her patience thinner — stepped forward and cut the ninth knot with a blade of bone. The blade had been sharpened on a stone from a river that had no name. The rival had carried it in her sleeve for three years. She had not expected to use it. She had hoped not to. Hope is not a rope.

The groom forgot his vows. Not slowly. Not gently. He stood at the altar with his mouth half open, his eyes the color of a winter river, his hands empty. He did not reach for the bride. He did not speak. He simply \textit{stopped} — like a timer whose owner has died, like a road that ends at a cliff, like a story whose last page has been torn out and burned.

The bride did not weep. Weeping would have been a kind of mercy, and she had none left. She stepped around him. She took his hand — cold, still, heavy as a bag of stones — and spoke for him.

\textit{``I give this ring. I take this name. I close this door. He is not here to refuse. That is not my problem. That is yours.''}

She said this to the priest. The priest was afraid. A priest who is afraid is a priest who has seen something he cannot explain to his god. He nodded. The bond held. The god did not object. The god, perhaps, was also afraid.

But the rival's hand withered where she had cut the cord — five fingers turned to three, the skin drawn tight as dried leather, the bone showing at the knuckles like the ribs of a wrecked ship. She wore a glove for the rest of her life. She never explained why. When children asked, she told them she had been bitten by a horse. The children did not believe her. But they stopped asking.

The bride kept the severed knot in a silver locket. The locket had belonged to her mother, who had died in childbirth, who had never tied any knots at all. She wore it for fifty years. She never opened it. She never needed to. The knot did not rattle. It did not whisper. It simply \textit{was} — a small weight against her chest, like a second heart that did not beat.

On her deathbed, she gave the locket to her daughter. The daughter was not the groom's child. The groom had never woken from his forgetting. He had lived another forty years, polite and vacant, a shell that ate and slept and walked to the window every morning to watch the road. He never said who he was waiting for. Perhaps he did not know.

\textit{``If you ever need to silence a liar,''} the bride said, \textit{``cut a knot. But be sure you know whose tongue will still speak afterwards. The blade does not ask whose blood it spills.''}

The daughter opened the locket. The knot was gone. Inside was a single tooth — the rival's, yellowed, cracked, with a thread of dried leather still tied around the root.

The daughter closed the locket. She did not ask. She had learned, from watching her mother, that some questions are doors. And some doors, once opened, cannot be closed again.

\section*{The Witch's Closing}

\textit{``The ninth knot is the thread that holds the story closed. Cut it, and the story unravels — but not always in the direction you expect. The blade remembers who swung it. The rope remembers who tied it. The silence remembers who was supposed to speak.''}

\section*{What Lurks in the Shadow of the Tale}

\subsection*{The Severed Knot}

A piece of wedding cord, kept in a locket or a box or a hollow tooth. It remembers the vow it once bound. It does not forget. It cannot. The knot is not a thing. It is a witness. When cut again — for it can always be cut again, and again, until nothing remains — it does not destroy the vow. It simply transfers the forgetting. Someone else forgets. Someone else pays. Someone else stands at the altar with empty hands and a mouth half open, and the priest nods because he is afraid, and the god does not object.

\paragraph{Tags / Mechanics:}
\begin{itemize}
\item [CURSE] [TRUTH] [OATH] [SILENCE]
\item As an action, cut the knot and speak the name of a vow you witnessed. The person who made that vow forgets they ever spoke it. They do not forget the person. They do not forget the place. They forget only the words. The shape of the forgetting is a circle: they will walk around the memory for the rest of their lives, never crossing the threshold, never knowing why they feel cold.
\item The cutter chooses who forgets. The cutter does not choose who pays. The GM chooses a cost: a withered hand (--1 die on physical rolls for one moon), a forgotten face (--1 die on social rolls for one moon), or a name traded (an enemy becomes a friend; a friend becomes a stranger; the cutter cannot remember which is which until the moon changes).
\item The knot can be cut once per season. It grows back smaller each time, tighter, more reluctant. The ninth cut unravels the cutter's own name. Not slowly. Not gently. Simply — \textit{stopped}.
\end{itemize}

\paragraph{Adventure Hook:}
A noble has been forced into a marriage contract she does not want. She possesses a ninth knot — her grandmother's, passed down through three widows, each of whom died with unopened eyes and a smile that did not reach their throats. Cutting the knot will make her betrothed forget his claim. But her grandmother's journal warns: \textit{``The third cut takes a finger. The sixth takes a memory. The ninth takes everything.''} The knot has been cut six times already. The thread is thin as a hair. The party must decide: cut again and accept the cost, find another way through the passes, or ask the question no one asks in the red tower: \textit{whose debt are we really paying?}

\section*{Colophon}
Red tower, dusk. The widow poured wine into eight cups. The ninth she emptied on the floor — not as a gift, not as an offering, but as a reminder. \textit{``For the ones who forgot,''} she said. \textit{``And for the ones who will. The floor remembers. The floor always remembers.''} The floor did not stain. The wine was older than the house, older than the bride, older than the knot. The wine had been waiting. It was still waiting.